% Notes-only C.49 candidate; full root/independent audit required.
\subsection{The actual framed images at levels twenty-four and sixty}

\begin{statement}{C.49}{Open}{No additional arithmetic-trivial character at the two requested frame levels}
Let $\mathcal H=\ker\lambda\subset\widehat{\mathrm{GT}}_0$.
For $g\in\mathcal H$ retain the actual cusp intertwiner and both
CM coordinates of C.12 and C.20:
\[
 V_g=\begin{pmatrix}1&t_g\\0&1\end{pmatrix},\qquad
 T_g=q(A_g)^{-1}V_g,\qquad W_g=q(\beta_g)^{-1}V_g.
                                                               \tag{C.49.1}
\]
For each coordinate separately fix the unique single global sign
making it equal to one modulo four. At $N=24,60$, the actual
image of $g\mapsto(t_g,T_g,W_g)\bmod N$ and the image of the
arithmetic subgroup $G_{\mathbb Q^{\rm ab}}$ are as follows:
\[
\begin{array}{c|c|c|c}
 N&\text{paired CM image}&\text{framed image on }\mathcal H
       &\text{arithmetic image}\\\hline
24&(C_2\times C_6)\times(C_2\times C_4)
   &C_2\times C_6\times C_2\times C_4&C_2^2\times C_{12}\\
60&C_6^2\times C_4^2
   &C_6^2\times C_4^2\times C_5&C_{12}\times C_6\times C_4\times C_5
\end{array}                                                    \tag{C.49.2}
\]
The corresponding framed/arithmetic orders are $96/48$ and
$2880/1440$. The arithmetic subgroup is exactly $\ker\Delta$
in each actual framed image. Hence its finite-order characters
trivial on the arithmetic image are precisely the trivial character
and $\Delta$, with values identified with $\{0,1/2\}\subset\mathbb Q/\mathbb Z$.
In particular the requested levels $2m$ with $m=12,30$ produce
no additional character of this kind from these retained frames.

An exact executable certificate computes the finite images and
has a residue-evaluation mode for every residue in them. These
are multiplier-one assertions. No order or presentation of the
arbitrary-multiplier framed image, no complete character basis
of the gentle group, and no smallest finite detector is asserted.
\end{statement}

\begin{proof}
\paragraph{Genuine signs and the local norm-one lift conditions.}
Use the signed matrices and quadratic-ring bases
\[
 S=\begin{pmatrix}0&1\\-1&0\end{pmatrix},\quad
 U=\begin{pmatrix}1&1\\0&1\end{pmatrix},\quad
 R=SU^{-1}=\begin{pmatrix}0&1\\-1&1\end{pmatrix},\qquad
 \omega=-R,\quad i=S.                                  \tag{C.49.3}
\]
Thus $\omega^2+\omega+1=0$ and $i^2=-1$; a coordinate
$a+b\omega$ is the matrix $aI-bR$, whereas $a+bi$ is $aI+bS$.
Their norms are $a^2-ab+b^2$ and $a^2+b^2$.
C.20 supplies the full actual adelic image
\[
 \{u\in\widehat{\mathbb Z[\omega]}^{\times}:u\bar u=1,
                      u\equiv1\pmod2\}/\{\pm1\}
 \ \times\
 \{u\in\widehat{\mathbb Z[i]}^{\times}:u\bar u=1,
                      u\equiv1\pmod2\}/\{\pm1\}.       \tag{C.49.4}
\]
The signs in these two quotients are independent of one another,
but each is a single sign shared by all coefficient primes.

For either ring, an actual norm-one element congruent to one
modulo two reduces modulo four to $1$ or $-1$.
Consequently every projective coordinate has a unique genuine
representative equal to one modulo four. This sign choice is
multiplicative. It must not be replaced by independent sign choices
at the odd primes.

Here are the exact local reductions at two. In the Eisenstein ring,
write $u=a+b\omega$ with $a$ odd and $b$ even.
The equation $a^2-ab+b^2\equiv1\pmod8$ forces $b\equiv0\pmod8$,
and permits every odd $a\pmod8$. All four residues lift to exact
norm-one units over $\mathbb Z_2$: for the chosen odd representative
$a\in\{1,3,5,7\}$, seek $u=a+8c\omega$. Its norm is one exactly when
\[
                  (a^2-1)/8-ac+8c^2=0.                 \tag{C.49.5}
\]
This polynomial has a root modulo two and odd derivative, so Hensel's
lemma gives an actual integral root. The full reduction modulo eight
is therefore $\{1,3,5,7\}\cong C_2^2$; modulo four it is
$\{\pm1\}$. After the global-sign normalization, the modulo-eight
part is $\{1,5\}\cong C_2$.

In the Gaussian ring, $u=a+bi\equiv1\pmod2$ has $a$ odd and
$b$ even. The exact norm equation implies $b\equiv0\pmod4$ and,
by reduction modulo sixteen, $a\equiv\pm1\pmod8$.
Every resulting residue is attained, using
\[
 \pm\frac{1+it}{1-it},\qquad t\in\{0,2\}\subset\mathbb Z_2.
                                                               \tag{C.49.6}
\]
Thus the genuine modulo-eight image is
$\{1,-1,1+4i,-1+4i\}\cong C_2^2$, and the normalized image is
$\{1,1+4i\}\cong C_2$.
The other norm-one residues modulo eight do not lift to genuine
norm-one units. In particular, simply enumerating the congruence
$a^2+b^2\equiv1\pmod8$ would overcount the image.

At an odd prime, every residue of norm one lifts to a genuine local
norm-one unit. Lift the residue arbitrarily with norm $n\equiv1\pmod p$
and multiply by the scalar square root of $n^{-1}$ in $1+p\mathbb Z_p$.
This works also at the ramified Eisenstein prime three, since two
is invertible. At each of three and five the Eisenstein reduction
is cyclic of order six, generated by $-\omega=R$; the Gaussian
reduction is cyclic of order four, generated by $i=S$.
These assertions are also checked by the complete finite ring tables
in the certificate.

Chinese remaindering now gives, in the normalized genuine signs,
\[
 H_{3,24}=C_2\times C_6,\quad H_{4,24}=C_2\times C_4,
 \qquad H_{3,60}=C_6\times C_6,\quad H_{4,60}=C_4\times C_4.
                                                               \tag{C.49.7}
\]
The full adelic surjectivity in C.20 proves that every such paired
residue is realized by an actual element of $\mathcal H$.
This is the full-group existence step; finite ring enumeration
alone would not prove it.

\paragraph{The actual shear restriction.}
Write the normalized modulo-three coordinates as
\[
                       T_3=R^k,\qquad W_3=S^l,
                 \quad k\bmod6,\quad l\bmod4.          \tag{C.49.8}
\]
The pure cubic frame $A_g$ lies in $F'$ for multiplier one.
This is the abelianization calculation of C.26: its order-three
primitive has abelianization zero when $\lambda=1$; changing the
pure marking and the inner normalization preserves that zero class.
Equivalently the invariant abelian vector under the order-three
peripheral permutation is zero, since its determinant-three
relation is injective over $\widehat{\mathbb Z}$.
Therefore the signed residue of
\[
                              A_g=V_g T_g^{-1}          \tag{C.49.9}
\]
belongs to the actual finite image of $F'$.

For transparency, the source generators in (C.49.3) are
$x=U^2$ and $y=SU^2S^{-1}$. Their finite signed images and
commutator images are computed from those two literal matrices,
with the following exact results:
\[
\begin{array}{c|r|r|c}
N&|q_N(F)|&|q_N(F')|&q_N(F')\text{ in CRT coordinates}\\\hline
24&768&16&\{I,5I\}\text{ at }8\ \times Q_8\text{ at }3\\
60&11520&960&\{I\}\text{ at }4\ \times Q_8\text{ at }3
                         \times\mathrm{SL}_2(\mathbb F_5)
\end{array}                                                    \tag{C.49.10}
\]
The signed pure image does not contain $-I$ modulo four, so these
same orders apply in the global projective quotient.
The certificate computes the commutator image as the subgroup
generated by the full conjugacy orbit of $[x,y]$; it checks the
projection orders and equality with the indicated direct products.

At eight the cubic coordinate is scalar, so (C.49.9) forces
$t_g\equiv0\pmod8$ when $N=24$. At four it forces
$t_g\equiv0\pmod4$ when $N=60$.
At three, the abelianization
$\mathrm{SL}_2(\mathbb F_3)/Q_8=C_3$ sends $S$ to one and
$R=SU^{-1}$ to the inverse class of $U$.
Equation (C.49.9) therefore gives $t_g+k\equiv0\pmod3$.
There is no additional condition at five, where the displayed
commutator image is the full group. Thus
\[
\begin{array}{c|l}
24&t_g\equiv0\pmod8,\quad t_g\equiv-k\pmod3,\\
60&t_g\equiv0\pmod4,\quad t_g\equiv-k\pmod3,
                         \quad t_g\bmod5\text{ unrestricted}.
\end{array}                                                    \tag{C.49.11}
\]
At twenty-four these congruences uniquely determine the shear from
$T$. Since all paired CM residues are actual, they already determine
the full framed image there.

\paragraph{The extra five-component is supplied by actual arithmetic.}
For level sixty, we must prove actual independent realization of the
last residue in (C.49.11); no finite compatibility claim replaces it.
The Tate normalization proved in C.20 makes the arithmetic shear
\[
                            t_\sigma=-8\kappa_2(\sigma).
                                                               \tag{C.49.12}
\]
Indeed, with Legendre parameter $z$ at the tangential cusp,
$q=z^2/256\,(1+z\mathbb Q[[z]])$. A compatible $n$-th root has
constant factor $2^{-8/n}$. The tangential splitting fixes the
formal $z^{1/n}$ and the formal unit root with constant one.
Thus Galois acts on the Tate quotient lift by the Kummer factor
$\zeta_n^{-8\kappa_2(\sigma)}$. The Tate submodule has cyclotomic
action and the quotient has action one, yielding exactly (C.49.12)
in the normalized lower-right-one representative. A global quadratic
twist disappears in the same projective normalization.
This is the actual Tate calculation of C.20, using
\cite[III, \S3, equations (23)--(26)]{PBC_MilneEC2021}.

The Kummer character $\kappa_2\bmod5$ maps
$G_{\mathbb Q^{\rm ab}}$ onto $C_5$.
For example $X^5-2$ is Eisenstein, its splitting field has group
$C_5\rtimes C_4$ with faithful action, and its derived subgroup
is $C_5$. The image of the closed derived subgroup of $G_{\mathbb Q}$
is that full derived subgroup. By Kronecker--Weber this closed
subgroup is $G_{\mathbb Q^{\rm ab}}$.
Since $-8$ is a unit modulo five, (C.49.12) is onto $C_5$ there.

C.20 and C.21 identify the arithmetic paired-CM image as
$\ker\Delta$. At level sixty it has order $576/2=288$, prime to five.
If the shear character modulo five vanished on the kernel of this
paired image, it would factor through a group of order $288$, which
is impossible for a nonzero character to $C_5$.
Consequently the arithmetic subgroup already realizes every shear
modulo five with both level-sixty CM coordinates equal to one.
Multiplying any actual CM realization by such an arithmetic element
proves that all residues allowed by (C.49.11) occur.
It proves at the same time that the arithmetic framed image has
exactly the extra $C_5$ factor asserted in (C.49.2).

\paragraph{The finite group laws and the arithmetic character.}
On $\mathcal H$, the three retained coordinates have the literal law
\[
 (t,T,W)(t',T',W')=(t+t',TT',WW').                       \tag{C.49.13}
\]
For the shear this follows from the unique normalized intertwiner;
for the two CM coordinates it is C.20's actual frame covariance.
The genuine sign normalizations modulo four are multiplicative.
At level sixty the shear kernel is $12\mathbb Z/60\cong C_5$,
and the choice $t\equiv0\pmod5$ splits it over the CM pair.
This gives the framed direct-product descriptions in (C.49.2).

The exact signed parity formula of C.47 is
\[
                         \Delta=k+l\pmod2.             \tag{C.49.14}
\]
It depends only on reduction modulo twelve, which divides both
levels. Thus reduction of C.20's adelic arithmetic kernel gives
exactly the finite parity kernel: every finite parity-zero CM pair
has an adelic lift, and every such lift has parity zero.
The preceding shear argument gives the same assertion for the full
framed data at sixty; at twenty-four the shear is already determined.
Inside $C_6\times C_4$, the subgroup of equal parities is the cyclic
order-twelve group generated by $(R,S)$. All other cyclic factors
are unconstrained. This proves the arithmetic direct products and
orders in (C.49.2).
The quotient of each actual framed image by its arithmetic image
is exactly $C_2$, so its finite-order arithmetic-trivial character
group is exactly the order-two group generated by (C.49.14).

\paragraph{Exact retained programs and limitations.}
The standalone program \texttt{code/pbc\_cm\_frames\_24\_60.py}
enumerates the signed pure groups of orders $768$ and $11520$,
their derived images, all actual local CM reduction classes and all
allowed shears. It checks the cyclic product generators and the
parity kernels, and verifies that the raw signed frame pair
$(A,\beta)$ is injectively recovered from the framed triples.
It also checks the local lift formulas at finite precision, while
the all-depth lifting proof is (C.49.5)--(C.49.6) and the odd-prime
square-root argument above. Its JSON and stdout are retained.

For direct evaluation, pass a level, shear, and the two coefficient
pairs in the bases of (C.49.3), for example
\[
 \texttt{--level 24 --shear 8 --T 9 8 --W 1 0}.
\]
The program returns $\Delta=1$. It checks membership in the proved
actual framed image, including the genuine two-adic lift condition,
and evaluates every member by (C.49.14). The standalone CLI checks
include arithmetic, nonarithmetic, bad-shear and nonlifting Gaussian
residues. Full profinite realization is justified by C.20 and the
actual arithmetic shear above, not by the finite enumeration itself.
The uncomputed arbitrary-multiplier image and the least finite
quotient in Task C.6 remain outside this statement.
\end{proof}
